\chapter*{前言}
\markboth{前言}{前言}
\addcontentsline{toc}{chapter}{前言}
% Source PDF page 011 (printed p. xi)
本书讨论凸优化。这是一类特殊的数学优化问题，包含最小二乘问题和线性规划问题。众所周知，最小二乘和线性规划有相当完备的理论，在众多领域中有应用，而且能够用数值方法高效求解。本书的基本观点是：对于范围更广的凸优化问题，同样可以这样说。

凸优化的数学研究已有大约一百年的历史，而近来几项相关进展又激发了人们对这一主题的兴趣。第一项进展是，人们认识到，20世纪80年代为求解线性规划而发展起来的内点法，也可以用于求解凸优化问题。这些新方法使某些新的凸优化问题，例如半正定规划和二阶锥规划，求解起来几乎与线性规划一样方便。

第二项进展是，人们发现，实际中出现的凸优化问题（除最小二乘和线性规划以外）比以往认为的更为普遍。自1990年以来，人们在自动控制系统、估计与信号处理、通信与网络、电子电路设计、数据分析与建模、统计和金融等领域发现了许多应用。凸优化在组合优化与全局优化中也得到了广泛应用，用于求最优值的界以及近似解。我们相信，还有许多凸优化的应用有待发现。

识别出一个问题是凸优化问题，或将它表述为凸优化问题，有很多好处。最基本的一点是，可以使用内点法或其他专门的凸优化方法，非常可靠、高效地求解它。这些求解方法足够可靠，可以嵌入计算机辅助设计或分析工具，甚至嵌入实时响应系统或自动控制系统。将问题表述为凸优化问题也有理论和概念上的好处。例如，与之相关的对偶问题，往往可以从原问题的角度作出有意义的解释，有时还能由此得到高效或分布式的求解方法。

我们认为，凸优化足够重要，凡是使用计算数学的人，都应该至少了解一些。在我们看来，在学习高等线性代数中的最小二乘、奇异值等内容以及线性规划之后，凸优化是一个自然的后续主题。

% Source PDF page 012 (printed p. xii)
\section*{本书的目标}
对于许多通用优化方法，通常的做法是直接把方法用于待求解的问题，看看效果如何。而要充分发挥凸优化的优势，就需要预先知道问题是凸的。当然，许多优化问题并不凸；识别其中的凸问题，或者把问题重新表述成凸问题，都可能不容易。

\begin{quote}
本书的主要目标是帮助读者掌握凸优化的实际运用，也就是获得识别、建立和求解凸优化问题所需要的技能与背景知识。
\end{quote}

掌握凸优化的实际运用，可能需要较多的数学准备，对于主要关心应用的读者尤其如此。根据我们的经验（主要来自电气工程和计算机科学专业的研究生），这份投入通常很有回报，有时回报非常可观。

有些关于线性规划和一般非线性规划的书着重讨论问题表述、建模与应用，另一些书则介绍凸优化理论，或内点法及其复杂度分析。本书的定位介于两者之间：它讨论一般凸优化，重点放在问题表述与建模上。

我们也应当说明本书的范围。本书的重点不是凸分析或凸优化的数学理论；已有一些教材很好地介绍了这些主题。本书也不对凸优化算法作全面综述。我们选择了少数几种好的算法，只介绍它们简单、理想化的形式（这些形式在实践中也确实很有效）。我们没有试图介绍求解凸问题的内点法或其他方法的最新研究进展。数值实现问题的讨论也作了较多简化，但我们认为，这些内容足以帮助有意使用这些方法的读者写出可用的实现；对于如何利用结构提高方法的效率，我们还作了较详细的介绍。对所述算法的复杂度理论，我们同样只作简化介绍。不过，我们会介绍自协调性这一重要概念，以及内点法复杂度分析的基本思想。

\section*{读者对象}
本书面向使用数学优化，或更广义地说，使用计算数学的研究人员、科学家和工程师。这自然包括直接从事优化和运筹学工作的人，也包括在计算机科学、经济学、金融、统计、数据挖掘以及许多科学与工程领域中使用优化的其他人。我们主要面向后一类读者，即凸优化的潜在使用者，而非人数相对较少的凸优化专家。

读者只需熟练掌握高等微积分和线性代数。如果还学过基本的数学分析（例如范数、收敛、初等拓扑）和初等概率论，应该就能读懂书中的每一处论证与讨论。我们也希望，
% Source PDF page 013 (printed p. xiii)
没有学过数学分析和概率论的读者，仍能掌握全部核心思想与要点。本书不要求读者事先接触过数值计算或优化，因为正文或附录会介绍所需的相关内容。

\section*{在课程中使用本书}
我们希望，本书能成为多种课程的主教材或辅助教材。自1995年以来，我们一直在斯坦福大学和加利福尼亚大学洛杉矶分校的线性、非线性及凸优化研究生课程中使用本书的草稿，这些课程都侧重工程应用。在一个学季的研究生课程中，我们能够讲授书中大部分内容，但无法逐一深入。一个学期的课程可以安排得更从容，介绍更多应用、更详细地讨论理论，并且可能安排一个简短的学生项目。连续两个学季的课程，则可以更充分地介绍线性规划、二次规划等基础主题（这些内容对重视应用的学生很有用），或安排一个规模更大的学生项目。

本书也可以作为传统线性与非线性优化课程的参考书或辅助教材，或用于包含部分凸优化内容的控制系统课程及其他应用领域课程。对于更侧重理论的凸优化课程，本书可以作为辅助教材，提供简单的实际例子。

\section*{致谢}
我们为本书准备材料已近十年。在这些年里，我们受益于许多人的反馈与建议，包括我们指导的研究生、选课学生，以及斯坦福大学、加利福尼亚大学洛杉矶分校和其他地方的同事。由于篇幅有限，而且记录不够完善，我们无法列出每一位作出贡献的人。不过，我们特别感谢 A.~Aggarwal、V.~Balakrishnan、A.~Bernard、B.~Bray、R.~Cottle、A.~d'Aspremont、J.~Dahl、J.~Dattorro、D.~Donoho、J.~Doyle、L.~El Ghaoui、P.~Glynn、M.~Grant、A.~Hansson、T.~Hastie、A.~Lewis、M.~Lobo、Z.-Q.~Luo、M.~Mesbahi、W.~Naylor、P.~Parrilo、I.~Pressman、R.~Tibshirani、B.~Van Roy、L.~Xiao 和 Y.~Ye。

J.~Jalden 和 A.~d'Aspremont 分别为第6.5.4节的时频分析例子和第6.5.5节的消费者偏好界例子作出了贡献。P.~Parrilo 提出了习题4.4和4.56。Igal Sason 对本书的细致阅读，为后续印次的修订提供了很大帮助。

我们还要特别感谢另外两位。Arkadi Nemirovski 激发了我们对凸优化最初的兴趣，并鼓励我们写作本书。我们也感谢 Kishan Baheti，他在本书的形成过程中起了关键作用。1994年，他鼓励我们向美国国家科学基金会申请一项将研究与课程建设相结合的项目，主题是凸优化及其工程应用。本书正是这一项目的直接成果，尽管问世的时间比预期晚了一些。

\begin{flushright}
Stephen Boyd\\
加利福尼亚州斯坦福\\[.8em]
Lieven Vandenberghe\\
加利福尼亚州洛杉矶
\end{flushright}
